% Part: sets-functions-relations
% Chapter: functions
% Section: inverses

\documentclass[../../../include/open-logic-section]{subfiles}

\begin{document}

\olfileid{sfr}{fun}{inv}
\olsection{Invers Fungsi}

\begin{explain}
Kita nganggep fungsi minangka pemetaan. Mula, pitakon kang lumrah
ngenani fungsi yaiku apa pemetaan iku bisa "diwalik". Upamane,
fungsi panerus $f(x) = x + 1$ bisa diwalik,
ing pangerten yen fungsi $g(y) = y - 1$ "mbalekake" asil saka $f$.

Nanging kita kudu ngati-ati. Sanajan definisi~$g$ nemtokake
fungsi $\Int \to \Int$, definisi iku ora nemtokake \emph{fungsi} $\Nat
\to \Nat$, amarga $g(0) \notin \Nat$. Dadi, sanajan ing kasus prasaja,
durung mesthi cetha apa sawijining fungsi bisa diwalik; iki bisa
gumantung marang domain lan kodomaine.

Pangerten invers fungsi ndadekake gagasan iki luwih tliti.
\end{explain}

\begin{defn}
Fungsi $g \colon B \to A$ iku \emph{invers} saka fungsi $f
\colon A \to B$ yen $f(g(y)) = y$ lan $g(f(x)) = x$ kanggo kabeh $x \in A$
lan $y \in B$.
\end{defn}

Yen $f$ duwe invers~$g$, kita kerep nulis $f^{-1}$ kanggo ngganti~$g$.

\begin{explain}
Saiki kita bakal nemtokake kapan fungsi duwe invers. Calon kang trep
kanggo invers saka $f\colon A \to B$ yaiku $g\colon B \to A$ kang
"ditetepake kanthi"
\[
g(y) = \text{$x$ "tunggal" kang nyukupi $f(x) = y$.}
\]
Nanging tandha petik ing "ditetepake kanthi" lan "tunggal" nuduhake
yen iki durung dadi definisi. Paling ora, cara iki ora mesthi lumaku
ing kabeh kasus. Supaya definisi iki nemtokake
fungsi, kudu ana siji lan mung siji~$x$ kang nyukupi $f(x) =
y$; asil saka~$g$ kudu ditemtokake kanthi tunggal. Kajaba iku, asil
kudu ditemtokake kanggo saben $y \in B$. Yen ana $x_1$ lan $x_2 \in
A$ kanthi $x_1 \neq x_2$ nanging $f(x_1) = f(x_2)$, mula $g(y)$ ora bakal
ditemtokake kanthi tunggal kanggo $y = f(x_1) = f(x_2)$. Lan yen ora ana~$x$
babar pisan kang nyukupi $f(x) = y$, mula $g(y)$ ora ditemtokake babar pisan.
Kanthi tembung liya, supaya $g$ bisa ditetepake, $f$~kudu !!{injective} lan
!!{surjective}.

Ayo dirembug sethithik mbaka sethithik. Pitakon mau dipara dadi loro:
yen diwenehi fungsi~$f\colon A \to B$, kapan ana fungsi $g\colon B \to A$
saengga $g(f(x)) = x$? Fungsi $g$ kang kaya mangkono "mbalekake" asil saka $f$,
lan diarani \emph{invers kiwa} saka~$f$. Kapindho, kapan ana
fungsi $h\colon B \to A$ saengga $f(h(y)) = y$? Fungsi $h$ kang kaya mangkono
diarani \emph{invers tengen} saka~$f$; $f$ "mbalekake" asil saka~$h$.
\end{explain}

\begin{prop}
Yen A ora kosong lan $f\colon A \to B$ iku !!{injective}, mula ana \emph{invers
kiwa}~$g\colon B \to A$ saka~$f$ saengga $g(f(x)) = x$ kanggo kabeh $x
\in A$.
Cathetan koreksi sumber OLFUN-001: sumber Inggris ora nyebut sarat
domain ora kosong kang digunakake ing bukti. Sarat umum kang pas yaiku
A ora kosong utawa B kosong; injeksi kosong menyang himpunan ora
kosong ora duwe invers kiwa.
\end{prop}

\begin{proof}
Upamakna A ora kosong lan $f\colon A \to B$ iku !!{injective}.
Gatekna sawijining $y \in B$.
Yen $y \in \ran{f}$, ana $x \in A$ saengga $f(x) = y$. Amarga
$f$ iku !!{injective}, mung ana siji~$x \in A$ kang kaya mangkono. Banjur kita bisa
netepake $g(y) = x$, yaiku $g(y)$ iku $x \in A$ "tunggal" kang nyukupi $f(x)
= y$. Yen $y \notin \ran{f}$, kita bisa metakake menyang sembarang~$a \in A$.
Dadi, kita bisa milih $a \in A$ lan netepake $g\colon B \to A$ kanthi:
\[
g(y) = \begin{cases}
    x & \text{yen $f(x) = y$}\\
    a & \text{yen $y \notin \ran{f}$.}
\end{cases}
\]
Fungsi iki ditetepake kanggo kabeh $y \in B$, amarga kanggo saben $y \in \ran{f}$
ana persis siji $x \in A$ kang nyukupi $f(x) = y$. Miturut definisi, yen
$y = f(x)$, mula $g(y) = x$, yaiku $g(f(x)) = x$.
\end{proof}

\begin{prob}
Tuduhna yen $f\colon A \to B$ duwe invers kiwa~$g$, mula $f$~iku
!!{injective}.
\end{prob}

\begin{prop}
    Yen $f\colon A \to B$ iku !!{surjective}, mula ana
    \emph{invers tengen}~$h\colon B \to A$ saka~$f$ saengga $f(h(y)) =
    y$ kanggo kabeh~$y \in B$.
\end{prop}

\begin{proof}
Upamakna $f\colon A \to B$ iku !!{surjective}. Gatekna sawijining $y \in
B$. Amarga $f$~iku !!{surjective}, ana $x_y \in A$ kanthi $f(x_y)
= y$. Banjur kita bisa netepake $h(y) = x_y$, yaiku kanggo saben $y \in B$
kita milih $x \in A$ saengga $f(x) = y$; amarga $f$~iku !!{surjective},
mesthi ana paling ora siji kang bisa dipilih.\footnote{Amarga $f$ iku
!!{surjective}, kanggo saben~$y \in B$ himpunan $\Setabs{x}{f(x) = y}$
ora kosong. Definisi~$h$ mbutuhake kita milih siji $x$
saka saben himpunan iki. Yen iki mesthi bisa ditindakake, sejatine ora
cetha; kamungkinan kanggo nggawe pilihan-pilihan iki dianggep minangka aksioma.
Kanthi tembung liya, proposisi iki nggunakake kang diarani Aksioma
Pilihan, prakara kang bakal \oliflabeldef{sth:choice::chap}{dirembug maneh ing \olref[sth][choice][]{chap}}{diliwati tanpa rembug jero}.
Nanging ing akeh kasus tartamtu, upamane nalika $A = \Nat$ utawa winates, utawa nalika $f$ iku !!{bijective},
Aksioma Pilihan ora dibutuhake. (Ing kasus mligi nalika $f$ iku !!{bijective}, kanggo saben $y \in B$ himpunan
$\Setabs{x}{f(x) = y}$ duwe persis siji !!{element}, mula ora perlu milih.)}
Miturut definisi, yen $x = h(y)$,
mula $f(x) = y$, yaiku kanggo saben $y \in B$, $f(h(y)) = y$.
\end{proof}

\begin{prob}
Tuduhna yen $f\colon A \to B$ duwe invers tengen~$h$, mula $f$~iku
!!{surjective}.
\end{prob}

\begin{explain}
  Kanthi nggabungake gagasan ing bukti sadurunge, saiki kita oleh asil yen saben
  !!{bijection} duwe invers, yaiku ana sawijining fungsi
  kang dadi invers kiwa lan uga invers tengen saka~$f$.
\end{explain}

\begin{prop}\ollabel{prop:bijection-inverse}
Yen $f\colon A \to B$ iku !!{bijective}, ana
fungsi~$f^{-1}\colon B \to A$ saengga kanggo kabeh $x \in A$,
$f^{-1}(f(x)) = x$ lan kanggo kabeh $y \in B$, $f(f^{-1}(y)) = y$.
\end{prop}

\begin{proof}
Minangka latihan.
\end{proof}

\begin{prob}
Buktekna \olref[sfr][fun][inv]{prop:bijection-inverse}. Kanggo iku, kudu
netepake~$f^{-1}$, nuduhake yen iku fungsi, lan nuduhake yen iku
invers saka~$f$, yaiku $f^{-1}(f(x)) = x$ lan $f(f^{-1}(y)) = y$ kanggo
kabeh $x \in A$ lan $y \in B$.
\end{prob}

\begin{explain}
Ana cara kang rada luwih umum kanggo nemokake invers. Ing
\olref[kin]{sec} wis kita deleng yen saben fungsi $f$ nuwuhake !!a{surjection} $f'
\colon A \to \ran{f}$ kanthi netepake $f'(x) = f(x)$ kanggo kabeh $x \in A$.
Cetha yen $f$~iku !!{injective}, mula $f'$~iku !!{bijective}, saengga
duwe invers tunggal miturut \olref{prop:bijection-inverse}. Kanthi
panganggone notasi kang rada longgar, kadhangkala invers saka $f'$
mung diarani "invers saka~$f$".
\end{explain}

\begin{prop}\ollabel{prop:left-right}%
  Tuduhna yen $f\colon A \to B$ duwe invers kiwa~$g$ lan invers
  tengen~$h$, mula $h = g$.
\end{prop}

\begin{proof}
  Minangka latihan.
\end{proof}

\begin{prob}
  Buktekna \olref[sfr][fun][inv]{prop:left-right}.
\end{prob}

\begin{prop}\ollabel{prop:inverse-unique}
Saben fungsi~$f$ duwe paling akeh siji invers.
\end{prop}

\begin{proof}
  Upamakna $g$ lan $h$ loro-lorone invers saka~$f$. Mula, mligine,
  $g$~iku invers kiwa saka~$f$ lan $h$~iku invers tengen. Miturut
  \olref{prop:left-right}, $g = h$.
\end{proof}

\end{document}
